% =====================================================================
%  ALL-k TWO-RADICAL KUMMER CLOSURE -- COMPLETED PROOF v1.0
%  Author: Will Lloyd (Lloyd Geometric Diagnostics)
%  Date: 18 July 2026
% =====================================================================
\documentclass[11pt,a4paper]{article}

\usepackage[margin=1.05in]{geometry}
\usepackage[T1]{fontenc}
\usepackage{lmodern}
\usepackage{microtype}
\usepackage{amsmath,amssymb,amsthm,mathtools}
\usepackage{booktabs}
\usepackage{array}
\usepackage{enumitem}
\usepackage{xcolor}
\usepackage[colorlinks=true,linkcolor=blue!60!black,citecolor=blue!60!black,
            urlcolor=blue!60!black]{hyperref}

\setlist[itemize]{nosep,leftmargin=*,topsep=2pt,partopsep=0pt}
\setlist[enumerate]{nosep,leftmargin=*,topsep=2pt,partopsep=0pt}
\widowpenalty=10000
\clubpenalty=10000
\brokenpenalty=10000

\newtheorem{theorem}{Theorem}[section]
\newtheorem{proposition}[theorem]{Proposition}
\newtheorem{lemma}[theorem]{Lemma}
\newtheorem{corollary}[theorem]{Corollary}
\theoremstyle{definition}
\newtheorem{definition}[theorem]{Definition}
\theoremstyle{remark}
\newtheorem{remark}[theorem]{Remark}

\newcommand{\Q}{\mathbb{Q}}
\newcommand{\Ftwo}{\mathbf F_2}
\newcommand{\Gal}{\operatorname{Gal}}
\newcommand{\Hom}{\operatorname{Hom}}
\newcommand{\Spec}{\operatorname{Spec}}
\newcommand{\ord}{\operatorname{ord}}
\newcommand{\wrp}{\mathop{\mathrm{wr}}}
\newcommand{\eps}{\varepsilon}
\newcommand{\gammam}{\gamma_k^{-}}

\title{\textbf{All-$k$ Two-Radical Kummer Closures}\\[4pt]
\large A canonical norm-derived radicand and maximal wreath monodromy\\[6pt]
\normalsize\fbox{Completed proof --- v1.0 --- 18 July 2026}}
\author{Will Lloyd\\
\small Lloyd Geometric Diagnostics Pty Ltd, Coffs Harbour, NSW, Australia}
\date{}

\begin{document}
\maketitle

\begin{abstract}
For the generic multiquadratic mass cover
\[
  w_i^2=u+N_i^2,\qquad 4M=\sum_{i=1}^k w_i,
\]
the base mass field has degree
\[
\delta_k=
\begin{cases}
2^{k-1},&k\text{ odd},\\
2^{k-1}-\frac12\binom{k}{k/2},&k\text{ even},
\end{cases}
\]
and generic normal-closure group $S_{\delta_k}$.  We construct, for every
$k\geq3$, a coefficient-free second radicand
\[
 \gamma_k=2\left(
 \sum_{r=0}^{\lfloor k/2\rfloor}u^r e_{k-2r}(w)
 +\prod_{i=1}^kN_i\right)
\]
that specializes at $k=4$ to the static horizon radicand.  Its companion
polynomial $\beta_k$ satisfies the uniform identity
\[
 \gamma_k(\gamma_k-4P)=4u\beta_k^2,
 \qquad P=\prod_iN_i.
\]
Even and odd signed-contact divisors give valuation rows $(1,0)$ and
$(1,1)$ for the channels $(u,\gamma_k)$.  We prove that all off-sheet
entries are even, so their Galois orbits form
$B\otimes I_{\delta_k}$ with
$B=\left(\begin{smallmatrix}1&0\\1&1\end{smallmatrix}\right)$.
Consequently all $2\delta_k$ conjugate square classes are independent.
The normal closure of
$K(\sqrt u,\sqrt{\gamma_k})/F$ therefore has the exact group
\[
 C_2^2\wrp S_{\delta_k}
\]
and order $2^{2\delta_k}\delta_k!$.  This proves the all-$k$ static
two-channel Kummer closure.  The theorem is algebraic for arbitrary $k$;
only the four-charge member is asserted here to carry the established
black-hole entropy interpretation.
\end{abstract}

\section{Introduction}

The four-charge horizon cover carries two quadratic channels above its
quintic mass field: the axial class $[u]$ and an entropy-sum class
$[\gamma_4]$.  Signed contact divisors separate these channels and yield the
static closure group $C_2^2\wrp S_5$ \cite{HorizonCover,KummerLift}.  It was
not previously clear whether the second class was forced by the general
multiquadratic norm structure or was an accident of the four-charge
truncation.

The issue is subtle only if ``second radicand'' is constrained.  An arbitrary
function field has many independent square classes, so searching for some
additional class would say nothing about the horizon construction.  The
candidate used here is instead forced by the two axial products
\[
 A_k^+=\prod_i(w_i+m),\qquad A_k^-=\prod_i(w_i-m),
 \qquad m^2=u.
\]
Their product is the square $P^2$, where $P=\prod_iN_i$.  Thus the
sum-compatible expression
\[
 A_k^++A_k^-+2P
\]
is the square of a sum after adjoining compatible square roots of $A_k^+$
and $A_k^-$.  It is invariant under $m\mapsto-m$, uses no adjustable
coefficients or selected auxiliary divisor, and recovers the published
four-charge normalization.  These properties pin the relevant algebraic
generalization.

Exact $k=3$ and $k=6$ falsifier computations first exposed the uniform
identity used below \cite{KillTest}.  The proof itself does not interpolate
from those cases.  It works for arbitrary $k\geq3$, proves the required local
valuation statements directly, and then invokes the already established
all-$k$ symmetric base monodromy \cite{WeightedMonodromy} and generic
reflection descent \cite{ReflectionDescent}.

\paragraph{Scope.}
The result is a theorem about the generic equal-weight multiquadratic mass
cover and its canonical norm-derived decoration.  At $k=4$ the decoration is
the physical static entropy-sum channel.  For other values of $k$, the theorem
asserts the algebraic closure of the displayed cover; it does not invent a
black-hole solution with an arbitrary number of independent charges.

\section{The generic mass cover}
\label{sec:setup}

Fix an integer $k\geq3$.  Let
\[
 F_0=\Q(N_1,\ldots,N_k),\qquad F=F_0(M),
\]
where the $N_i$ and $M$ are algebraically independent.  On the
multiquadratic curve over $F_0$ put
\begin{equation}
 w_i^2=u+N_i^2\quad(1\leq i\leq k),
 \qquad \sigma=\sum_{i=1}^kw_i.
 \label{eq:curve}
\end{equation}
The mass fibre is $\sigma=4M$.  Let $\mathcal N_k(u;M)$ denote the irreducible
mass polynomial, let $u$ be one of its roots, and put
\[
 K=F(u).
\]

We use two established results.  They are stated here in precisely the form
needed by the new argument.

\begin{theorem}[All-$k$ base monodromy \cite{WeightedMonodromy}]
\label{thm:base}
The extension $K/F$ is generically separable of degree
\begin{equation}
 \delta_k=
 \begin{cases}
 2^{k-1},&k\text{ odd},\\[2mm]
 2^{k-1}-\dfrac12\binom{k}{k/2},&k\text{ even}.
 \end{cases}
 \label{eq:delta}
\end{equation}
Its normal closure $E/F$ has Galois group
\[
 \Gal(E/F)\cong S_{\delta_k},
\]
acting in the natural way on the $\delta_k$ embeddings of $K$ into $E$.
\end{theorem}

\begin{theorem}[Reflection descent \cite{ReflectionDescent}]
\label{thm:descent}
Let $m^2=u$ and let $W$ be the generic signed mass field containing
$m,w_1,\ldots,w_k$.  The involution
\[
 \tau(m)=-m,\qquad \tau(w_i)=w_i
\]
has fixed field $K$.  In particular,
\[
 w_i\in K\qquad(1\leq i\leq k).
\]
\end{theorem}

The descent statement is crucial: the symmetric expressions constructed
below are not merely functions on an auxiliary radical sheet.  They belong to
the degree-$\delta_k$ mass field itself.

\section{The canonical second radicand}
\label{sec:radicand}

Write $e_j(w)=e_j(w_1,\ldots,w_k)$ for the $j$th elementary symmetric
polynomial, with $e_0(w)=1$, and put
\begin{align}
 \alpha_k
 &=\sum_{r=0}^{\lfloor k/2\rfloor}u^r e_{k-2r}(w),
 \label{eq:alpha}\\
 \beta_k
 &=\sum_{r=0}^{\lfloor(k-1)/2\rfloor}u^r e_{k-1-2r}(w).
 \label{eq:beta}
\end{align}
By Theorem~\ref{thm:descent}, both elements lie in $K$.

\begin{lemma}[Parity decomposition]
\label{lem:decomposition}
In $K(m)$ one has
\begin{equation}
 A_k^+:=\prod_{i=1}^k(w_i+m)=\alpha_k+m\beta_k,
 \qquad
 A_k^-:=\prod_{i=1}^k(w_i-m)=\alpha_k-m\beta_k.
 \label{eq:Apm}
\end{equation}
\end{lemma}

\begin{proof}
Expanding the first product by elementary symmetric functions gives
\[
 \prod_{i=1}^k(w_i+m)
 =\sum_{j=0}^k m^{k-j}e_j(w).
\]
The terms with $k-j$ even are exactly
\[
 u^{(k-j)/2}e_j(w),\qquad j\equiv k\pmod2,
\]
and their sum is $\alpha_k$.  The remaining terms contain one factor of $m$;
after removing it, their sum is $\beta_k$.  Replacing $m$ by $-m$ fixes the
even part and negates the odd part, proving both identities.
\end{proof}

\begin{proposition}[Uniform norm identity]
\label{prop:normidentity}
Let
\[
 P=\prod_{i=1}^kN_i.
\]
Then
\begin{equation}
 \alpha_k^2-u\beta_k^2=P^2.
 \label{eq:pell}
\end{equation}
Consequently, for
\begin{equation}
 \gamma_k=2(\alpha_k+P),
 \label{eq:gamma}
\end{equation}
one has
\begin{equation}
 \boxed{\gamma_k(\gamma_k-4P)=4u\beta_k^2.}
 \label{eq:masteridentity}
\end{equation}
\end{proposition}

\begin{proof}
Multiplying the two products in \eqref{eq:Apm} gives
\[
 A_k^+A_k^-
 =\prod_{i=1}^k(w_i^2-m^2)
 =\prod_{i=1}^kN_i^2
 =P^2,
\]
because $w_i^2-m^2=(u+N_i^2)-u=N_i^2$.  The left-hand side is
$\alpha_k^2-u\beta_k^2$, proving \eqref{eq:pell}.  Finally,
\[
 \gamma_k(\gamma_k-4P)
 =4(\alpha_k+P)(\alpha_k-P)
 =4(\alpha_k^2-P^2)
 =4u\beta_k^2.
\]
\end{proof}

\begin{proposition}[Canonicity and exhaustion of the cross-term choices]
\label{prop:canonicity}
Choose square roots $x_\pm^2=A_k^\pm$ compatibly so that $x_+x_-=P$.
Then
\[
 \gamma_k=(x_++x_-)^2.
\]
The opposite cross-term
\[
 \gammam:=A_k^++A_k^--2P=2(\alpha_k-P)=\gamma_k-4P
\]
satisfies, in $K^*/K^{*2}$,
\begin{equation}
 [\gammam]=[u][\gamma_k].
 \label{eq:minusclass}
\end{equation}
Thus the two sum/difference choices generate no square class beyond
$[u]$ and $[\gamma_k]$.
\end{proposition}

\begin{proof}
The first assertion follows from
\[
 (x_++x_-)^2=A_k^++A_k^-+2P=2(\alpha_k+P).
\]
Equation~\eqref{eq:masteridentity} reads
$\gamma_k\gammam=4u\beta_k^2$.  Since $4\beta_k^2$ is a square, taking
square classes gives \eqref{eq:minusclass} wherever the displayed elements
are nonzero.  Since $\gamma_k$ is a unit at the even contact constructed
below, it is not the zero element of $K$; the same identity then shows that
the generic square-class statement is well defined.
\end{proof}

\begin{remark}
For $k=4$, equations \eqref{eq:alpha}--\eqref{eq:gamma} give
\[
 \alpha_4=e_4(w)+ue_2(w)+u^2,\qquad
 \beta_4=e_3(w)+ue_1(w),\qquad
 \gamma_4=2\bigl(e_4(w)+ue_2(w)+u^2+P\bigr),
\]
which is exactly the globally normalized static R9 radicand.  The shorter
expression $2(P+e_4(w))$ is only its specialization at $u=0$.
\end{remark}

\section{Signed-contact divisors}
\label{sec:contacts}

For a sign vector $\eps=(\eps_1,\ldots,\eps_k)\in\{\pm1\}^k$, define
\begin{equation}
 \ell_\eps=4M-\sum_{i=1}^k\eps_iN_i,
 \qquad s_\eps=\prod_{i=1}^k\eps_i.
 \label{eq:wall}
\end{equation}
The base divisor $D_\eps$ is the generic point of $\ell_\eps=0$, localized
away from
\begin{equation}
 P\prod_{\eta\neq\eps}\ell_\eta
 \left(\sum_{i=1}^k\frac{\eps_i}{N_i}\right)=0.
 \label{eq:localization}
\end{equation}
The distinct linear forms $\ell_\eta$ define distinct divisors.  The final
factor in \eqref{eq:localization} is not the zero rational function: after
multiplication by $P$ it is a sum of distinct monomials
$\eps_i\prod_{j\neq i}N_j$.  Hence this localization is nonempty for every
sign vector.

\begin{lemma}[Simple signed contact]
\label{lem:simplecontact}
Above $D_\eps$ there is a unique mass sheet on which
\begin{equation}
 u=0,\qquad w_i=\eps_iN_i.
 \label{eq:contactpoint}
\end{equation}
On that sheet $u$ is a uniformizer.  Equivalently, its valuation satisfies
$v_\eps(u)=1$.
\end{lemma}

\begin{proof}
On the branch through \eqref{eq:contactpoint}, the binomial expansion gives
\[
 w_i
 =\eps_iN_i\left(1+\frac{u}{N_i^2}\right)^{1/2}
 =\eps_iN_i+\frac{\eps_i}{2N_i}u+O(u^2).
\]
Substitution in $4M=\sum_iw_i$ yields
\begin{equation}
 \ell_\eps
 =\frac{u}{2}\sum_{i=1}^k\frac{\eps_i}{N_i}+O(u^2).
 \label{eq:localexpansion}
\end{equation}
The coefficient of $u$ is a unit by \eqref{eq:localization}; hence the local
implicit-function theorem, or equivalently formal inversion of
\eqref{eq:localexpansion}, shows that $u$ and $\ell_\eps$ have the same order.
Thus $v_\eps(u)=1$.

At $u=0$, every signed radical branch has mass value
$\sum_i\eta_iN_i$.  On $D_\eps$, the other values differ by the units
$\ell_\eta$.  Therefore no other mass sheet has $u=0$ at the generic point of
$D_\eps$, proving uniqueness.
\end{proof}

\begin{lemma}[Contact values of $\alpha_k$ and $\beta_k$]
\label{lem:contactvalues}
At the contact \eqref{eq:contactpoint},
\begin{equation}
 \alpha_k(0)=s_\eps P,\qquad
 \beta_k(0)=s_\eps P\sum_{i=1}^k\frac{\eps_i}{N_i}.
 \label{eq:contactvalues}
\end{equation}
In particular, $\beta_k$ is a unit in the localization
\eqref{eq:localization}.
\end{lemma}

\begin{proof}
Every term in \eqref{eq:alpha} except $e_k(w)$ is divisible by $u$, and every
term in \eqref{eq:beta} except $e_{k-1}(w)$ is divisible by $u$.  Hence
\[
 \alpha_k(0)=\prod_i\eps_iN_i=s_\eps P.
\]
Moreover,
\begin{align*}
 \beta_k(0)
 &=e_{k-1}(\eps_1N_1,\ldots,\eps_kN_k)\\
 &=\sum_{i=1}^k\prod_{j\neq i}\eps_jN_j
 =s_\eps P\sum_{i=1}^k\frac{\eps_i}{N_i}.
\end{align*}
The claimed unit statement follows from \eqref{eq:localization}.
\end{proof}

\begin{proposition}[The two universal parity rows]
\label{prop:rows}
At an even signed contact, $s_\eps=+1$, one has
\[
 \bigl(v_\eps(u),v_\eps(\gamma_k)\bigr)\equiv(1,0)\pmod2.
\]
At an odd signed contact, $s_\eps=-1$, one has
\[
 \bigl(v_\eps(u),v_\eps(\gamma_k)\bigr)\equiv(1,1)\pmod2.
\]
Thus the sheet-level valuation matrix is
\begin{equation}
 B=
 \begin{pmatrix}
 1&0\\
 1&1
 \end{pmatrix}\in\operatorname{GL}_2(\Ftwo).
 \label{eq:B}
\end{equation}
\end{proposition}

\begin{proof}
Lemma~\ref{lem:simplecontact} gives $v_\eps(u)=1$ in both cases.  If
$s_\eps=+1$, then Lemma~\ref{lem:contactvalues} gives
\[
 \gamma_k(0)=2(P+P)=4P,
\]
which is a unit; hence $v_\eps(\gamma_k)=0$.

If $s_\eps=-1$, then $\gamma_k(0)=0$, while
$\gamma_k(0)-4P=-4P$ is a unit.  Taking valuations in
\eqref{eq:masteridentity}, and using that $\beta_k$ is a unit, gives
\[
 v_\eps(\gamma_k)
 =v_\eps(u)+2v_\eps(\beta_k)=1.
\]
These are the two rows of \eqref{eq:B}; its determinant is one in $\Ftwo$.
\end{proof}

\section{Off-sheet parity and orbit saturation}
\label{sec:orbit}

Let $E/F$ be the normal closure of $K/F$ and let
\[
 I=\Hom_F(K,E),\qquad |I|=\delta_k.
\]
For $i\in I$, write
\[
 u_i=i(u),\qquad \gamma_i=i(\gamma_k),\qquad \beta_i=i(\beta_k).
\]
Choose a prime of $E$ above the even or odd divisor $D_\eps$ and above the
distinguished contact branch of Lemma~\ref{lem:simplecontact}.  The next lemma
is the point that prevents hidden conjugate entries from corrupting the local
matrix.

\begin{lemma}[Off-sheet entries are even]
\label{lem:offsheet}
At such a prime, every off-sheet axial element $u_j$ is a unit.  Every
off-sheet entropy element $\gamma_j$ has even valuation.  Hence all
off-sheet entries in the valuation-parity row are zero.
\end{lemma}

\begin{proof}
The unique-contact statement in Lemma~\ref{lem:simplecontact} says that only
the distinguished mass root vanishes at the generic point of $D_\eps$.
Therefore $v(u_j)=0$ for every off-sheet index $j$.

All $u_j,w_{r,j}$ and $\gamma_j$ are regular after the localization
\eqref{eq:localization}: the mass polynomial has nonzero leading coefficient
there, each $w_{r,j}$ is integral over the corresponding local ring, and
$\gamma_j$ is polynomial in these elements.  Thus $v(\gamma_j)\geq0$.
If $v(\gamma_j)=0$, there is nothing to prove.  If it is positive, then
$P$ being a unit implies that $\gamma_j-4P$ is a unit.  Applying the conjugate
of \eqref{eq:masteridentity} gives
\[
 v(\gamma_j)
 =v(u_j)+2v(\beta_j)=2v(\beta_j),
\]
which is even.
\end{proof}

\begin{proposition}[Full conjugate parity matrix]
\label{prop:fullmatrix}
There exist $2\delta_k$ divisorial valuations of $E$ whose parity matrix on
the ordered conjugate radicands
\[
 (u_i)_{i\in I}\mid(\gamma_i)_{i\in I}
\]
is, after ordering rows and columns,
\begin{equation}
 \mathcal V=B\otimes I_{\delta_k}
 =
 \begin{pmatrix}
 I_{\delta_k}&0\\
 I_{\delta_k}&I_{\delta_k}
 \end{pmatrix}.
 \label{eq:fullmatrix}
\end{equation}
In particular,
\[
 \operatorname{rank}_{\Ftwo}(\mathcal V)=2\delta_k.
\]
\end{proposition}

\begin{proof}
Choose one even sign vector and one odd sign vector.  Proposition~\ref{prop:rows}
and Lemma~\ref{lem:offsheet} produce the two rows of $B$, supported on one
distinguished embedding.  By Theorem~\ref{thm:base}, the group
$\Gal(E/F)=S_{\delta_k}$ acts transitively on $I$.  Acting on the selected primes transports each row to
every embedding without mixing the two channel coordinates.  The resulting
matrix is exactly \eqref{eq:fullmatrix}.  Since $B$ and
$I_{\delta_k}$ are invertible over $\Ftwo$, so is their Kronecker product.
\end{proof}

\section{The all-\texorpdfstring{$k$}{k} Kummer closure}
\label{sec:closure}

Define the decorated sheet field
\begin{equation}
 H_k=K(\sqrt u,\sqrt{\gamma_k})
 \label{eq:H}
\end{equation}
and its conjugate radical composite
\begin{equation}
 L_k=E\bigl(\sqrt{u_i},\sqrt{\gamma_i}:i\in I\bigr).
 \label{eq:L}
\end{equation}
The field $L_k$ is the normal closure of $H_k/F$: every $F$-embedding sends
each displayed square root to a signed square root of another displayed
radicand, while $L_k$ contains all such conjugates.

\begin{lemma}[Valuation rank detects all square classes]
\label{lem:kummerrank}
The $2\delta_k$ classes
\[
 [u_i],\ [\gamma_i]\in E^*/E^{*2}\qquad(i\in I)
\]
are linearly independent over $\Ftwo$.
\end{lemma}

\begin{proof}
Suppose a product
\[
 \prod_{i\in I}u_i^{a_i}\gamma_i^{b_i}
\]
is a square in $E$, with $a_i,b_i\in\Ftwo$.  Evaluating its valuation modulo
two at each of the $2\delta_k$ primes in
Proposition~\ref{prop:fullmatrix} says
\[
 \mathcal V(a_i,b_i)_{i\in I}^{\mathsf T}=0.
\]
The matrix $\mathcal V$ is invertible, so every $a_i$ and $b_i$ is zero.
\end{proof}

\begin{theorem}[All-$k$ two-radical Kummer closure]
\label{thm:main}
For every $k\geq3$, on the generic locus described above,
\begin{align}
 [H_k:K]&=4, & [H_k:F]&=4\delta_k,\label{eq:sheetdegree}\\
 \Gal(L_k/E)&\cong(C_2^2)^{\delta_k},
 &
 \Gal(L_k/F)&\cong C_2^2\wrp S_{\delta_k}.
 \label{eq:groups}
\end{align}
Consequently,
\begin{equation}
 |\Gal(L_k/F)|=2^{2\delta_k}\delta_k!.
 \label{eq:order}
\end{equation}
Equivalently, the conjugate Kummer relation module is zero.
\end{theorem}

\begin{proof}
Lemma~\ref{lem:kummerrank} proves that the span in $E^*/E^{*2}$ of the
$2\delta_k$ conjugate radicands has full dimension $2\delta_k$.  Kummer
theory in characteristic different from two therefore gives
\[
 \Gal(L_k/E)\cong(C_2)^{2\delta_k}\cong(C_2^2)^{\delta_k}.
\]
In particular, the two classes $[u]$ and $[\gamma_k]$ are independent in
$K^*/K^{*2}$, proving $[H_k:K]=4$ and hence \eqref{eq:sheetdegree}.

Restriction to $E$ gives an exact sequence
\[
 1\longrightarrow(C_2^2)^{\delta_k}
 \longrightarrow\Gal(L_k/F)
 \longrightarrow S_{\delta_k}\longrightarrow1.
\]
Every automorphism permutes the embedding labels and may change the signs of
the two square roots above each label.  Thus $\Gal(L_k/F)$ embeds in the
permutational wreath product
\[
 (C_2^2)^{\delta_k}\rtimes S_{\delta_k}.
\]
Its kernel is the full base group $(C_2^2)^{\delta_k}$ and its projection to
$S_{\delta_k}$ is surjective.  A subgroup of a semidirect product that
contains the full normal factor and surjects onto the quotient is the whole
semidirect product.  This proves \eqref{eq:groups}; taking orders gives
\eqref{eq:order}.  The relation module is the kernel of the map from the free
$2\delta_k$-dimensional permutation module to the conjugate square-class
span, and Lemma~\ref{lem:kummerrank} shows that kernel is zero.
\end{proof}

\begin{corollary}[No additional cross-term channel]
\label{cor:nothird}
Adjoining $\sqrt{\gammam}$ to $H_k$ does not enlarge it:
\[
 K(\sqrt u,\sqrt{\gamma_k},\sqrt{\gammam})
 =K(\sqrt u,\sqrt{\gamma_k}).
\]
\end{corollary}

\begin{proof}
Equation~\eqref{eq:masteridentity} gives
\[
 \sqrt{\gammam}=\frac{2\beta_k\sqrt u}{\sqrt{\gamma_k}}
\]
after a compatible choice of signs.
\end{proof}

\section{Low-\texorpdfstring{$k$}{k} specializations}
\label{sec:specializations}

The first four nontrivial members are shown in Table~\ref{tab:lowk}.  The
$k=3,5,6$ contact calculations are exact checks of the uniform proof; the
$k=4$ row recovers the established static horizon crown.

\begin{table}[htbp]
\centering
\caption{First all-$k$ two-channel closures.}
\label{tab:lowk}
\renewcommand{\arraystretch}{1.18}
\begin{tabular}{@{}c c c c l@{}}
\toprule
$k$ & $\delta_k$ & Kummer rank & $[H_k:F]$ & Normal-closure group\\
\midrule
3 & 4  & 8  & 16 & $C_2^2\wrp S_4$\\
4 & 5  & 10 & 20 & $C_2^2\wrp S_5$\\
5 & 16 & 32 & 64 & $C_2^2\wrp S_{16}$\\
6 & 22 & 44 & 88 & $C_2^2\wrp S_{22}$\\
\bottomrule
\end{tabular}
\end{table}

For reference, the explicit symmetric expressions are
\begin{align*}
 \alpha_3&=e_3+ue_1,
 &\beta_3&=e_2+u,\\
 \alpha_4&=e_4+ue_2+u^2,
 &\beta_4&=e_3+ue_1,\\
 \alpha_5&=e_5+ue_3+u^2e_1,
 &\beta_5&=e_4+ue_2+u^2,\\
 \alpha_6&=e_6+ue_4+u^2e_2+u^3,\\
 \beta_6&=e_5+ue_3+u^2e_1.
\end{align*}
The alternating pattern is not empirical: it is exactly the parity
decomposition proved in Lemma~\ref{lem:decomposition}.

\section{Exceptional loci and theorem boundary}
\label{sec:exceptions}

The theorem is generic, and its proof identifies the excluded set rather
than hiding it in a single unspecified discriminant.  The argument localizes
away from the following divisors and their higher intersections:

\begin{enumerate}
\item $P=0$, where one or more charge parameters vanish and the norm square
      $P^2$ degenerates;
\item $N_i^2=N_j^2$, where the multiquadratic degree or normalized sheet
      structure can drop;
\item the reciprocal sub-balance loci
      $\sum_i\eps_i/N_i=0$, where the chosen signed contact ceases to be
      simple or $\beta_k$ ceases to be a unit;
\item intersections $\ell_\eps=\ell_\eta=0$ of distinct signed walls, where
      more than one contact root can occur simultaneously;
\item the base discriminant and leading-coefficient loci of the mass
      polynomial; and
\item characteristic two, where valuation parity and ordinary quadratic
      Kummer theory require replacement.
\end{enumerate}

Each exclusion is proper.  In particular, the reciprocal sub-balance
function is not identically zero, and the signed wall forms are pairwise
distinct.  Thus the open locus used by the proof is nonempty.  Because the
base group and Kummer rank are generic open conditions, the intersection of
their nonempty open loci is again nonempty in the irreducible parameter
space.

For even $k$, balanced sign classes alter the pole count at infinity and are
responsible for the binomial correction in $\delta_k$.  They do not affect the
finite signed-contact calculation: the local expansion
\eqref{eq:localexpansion}, the matrix $B$, and the off-sheet identity remain
unchanged.

\section{Consequences}

The theorem closes the static all-$k$ Kummer relation-module problem for the
canonical norm-derived two-channel cover.  Once the symmetric base group is
known, no new large Galois computation is required: the full result is forced
by two sheet-local contact types and the global identity
\eqref{eq:masteridentity}.

The formula also explains why the $k=4$ calculation was reusable.  The
published relation
\[
 [\gamma_4-4P]=[u][\gamma_4]
\]
is not a special multiplicative coincidence.  It is the $k=4$ member of
\[
 [\gamma_k-4P]=[u][\gamma_k]
\]
for every $k$.  What is special at four charges is the established physical
entropy interpretation, not the underlying Kummer algebra.

The same architecture can be applied to further decorated channels.  For any
additional canonical radicand, the remaining work is sharply constrained:
identify its divisor types, prove regularity, compute the new sheet-level
parity rows, and test whether the enlarged matrix remains invertible.  The
base permutation group and the two channels proved here need not be rebuilt.

\section*{Reproducibility note}

The exact $k=3$ and $k=6$ symbolic identities, contact mass norms, simple-root
checks and parity matrices are recorded in the machine-verifiable evidence
package \cite{KillTest}.  The $k=5$ specialization was independently replayed
with the same exact builder, giving degree $16$, contact rows $(1,0)$ and
$(1,1)$, and Kummer rank $32$.  These calculations are checks of the theorem,
not premises of its proof.

{\small
\begin{thebibliography}{9}

\bibitem{WeightedMonodromy}
W.~Lloyd,
\emph{Generic Symmetric Monodromy of Weighted Multiquadratic Sums: A Complete
All-$k$ Theorem with a Kummer--Wreath Lifting Criterion},
completed theorem-and-proof paper, v1.0, 16 July 2026.

\bibitem{ReflectionDescent}
W.~Lloyd,
\emph{Audit Report---R6 + R7 Generic Birationality and Reflection Descent},
proof audit, 2026.

\bibitem{KummerLift}
W.~Lloyd,
\emph{Conjugate Kummer Modules and Wreath-Product Closures},
general theorem and proof, 10 July 2026.

\bibitem{HorizonCover}
W.~Lloyd,
\emph{Galois Theory of the Horizon Cover: Abel--Ruffini Obstruction,
Kummer Lifts, and Wreath-Product Monodromy at Four Charges},
technical preprint, v1.0, 11 July 2026.

\bibitem{KillTest}
W.~Lloyd and OpenAI Codex,
\emph{The $k=3/k=6$ Radicand Kill Test},
exact machine-verifiable evidence report, v1.0, 18 July 2026.

\end{thebibliography}
}

\end{document}
